\section{Shared Resources}

\subsection{Race Conditions}

\begin{definition}
\textbf{Race condition: } correctness depends on specific interleaving.
\end{definition}

\begin{definition}
\textbf{Low-level race condition: } insufficient synchronization of memory access (Data Races).
\textbf{High-level race condition: } wrong decomposition/algorithm (bad interleavings despite mutual exclusion).
\end{definition}

Generally we want one of three things when it comes to shared ressources:
\begin{itemize}
    \item \textbf{Thread local (Isolated Mutability):} We dont use the ressource of other threads.
    \item \textbf{Read only (Immutable):} We only read and do not change the ressource
    \item \textbf{Mutual exclusion (Synchronization):} We use the ressource of other threads but we ensure that only one thread accesses it at a time.
\end{itemize}

\subsection{Guidelines for Concurrent programming}

\begin{important}
\textbf{Thread-Locality:} Only share ressources if needed.
\end{important}

\begin{important}
\textbf{Immutability:} Use immutable data structures whenever possible.
\end{important}

\begin{important}
\textbf{Consitent locking:} Use locks to ensure mutual exclusion.
\end{important}

\begin{definition}
\textbf{Lock granularity } refers to the amount of code that is proteted by a lock in ways:
\begin{itemize}
        \item Size of synchronized blocks
        \item Number of objects a lock is used for
\end{itemize}
\end{definition}

\begin{theorem}
\textbf{Coarse-grained locking} is easier to implement but leads to less concurrency. \textbf{Fine-grained locking} is harder to implement without creating bugs but leads to more concurrency.
\end{theorem}

\begin{important}
Start with coarse-grained (simpler) and move to fine-
grained (performance) only if contention on the coarser locks becomes an issue.
\end{important}

\begin{definition}
\textbf{Atomicity: } An operation is atomic if it is performed as a single, indivisible unit. 
\end{definition}

\begin{important}
    Do not do expensive computations or I/O in critical
    sections, but also don’t introduce race conditions
\end{important}